\chapter*{Résumé}
\addcontentsline{toc}{chapter}{Résumé}

  Ce mémoire porte sur Smart Focus, un système que nous avons conçu et développé pour analyser en continu l'état de l'utilisateur pendant ses sessions de
  travail ou de révision : concentration, posture, fatigue, stress, à partir d'une caméra et de capteurs intégrés dans un boîtier basé sur Raspberry Pi.

  L'idée de départ était simple : plutôt que de laisser l'utilisateur estimer lui-même son niveau d'attention, le système le mesure objectivement et réagit en
  conséquence. Une application mobile complète le dispositif avec des fonctionnalités de suivi des sessions, de planification intelligente, et d'un chatbot
  pédagogique basé sur une approche \ac{RAG}, capable de répondre à des questions sur des documents uploadés par l'utilisateur.

  Le développement s'est déroulé en trois sprints Agile, chacun livrant un incrément fonctionnel testé sur du matériel réel : authentification et profil,
  module de vision et calcul du score, puis les fonctionnalités pédagogiques.

  Au final, Smart Focus est un système complet, du traitement temps réel sur Raspberry Pi jusqu'à l'interface mobile, pensé pour s'intégrer dans la routine
  quotidienne sans friction.

\bigskip
\noindent
\textbf{Mots-clés :} IoT, intelligence artificielle, concentration, stress, posture, fatigue, Raspberry Pi, RAG.

\cleardoublepage

\chapter*{Abstract}
\addcontentsline{toc}{chapter}{Abstract}

  Smart Focus is an IoT and AI system that monitors a user's concentration, posture, fatigue, and stress in real time during work or study sessions. The core
  hardware is a Raspberry Pi enclosure fitted with a camera and sensors; a companion mobile application handles session tracking, intelligent scheduling, and
  learning support through a RAG-based chatbot.

  The system was built around a straightforward question: instead of asking users to estimate their own attention level, can we measure it objectively and give
   them actionable feedback? The vision pipeline runs continuously on the Raspberry Pi, extracting behavioral cues: gaze direction, posture, facial signals,
  and combining them into a live concentration score.

  Development followed an Agile methodology across three sprints, each delivering working functionality tested on real hardware before moving forward. The
  architecture connects an embedded processing layer, a Python backend, and a Flutter mobile interface, three components that had to work together in real
  time.

\bigskip
\noindent
\textbf{Keywords:} IoT, artificial intelligence, concentration, stress, posture, fatigue, Raspberry Pi, RAG.